% Separately instantiated corrected French for printed p.121.
% Decision: TOH-P121-R1-SRC-DISPOSITION-001.
% No authorial error is established; historical wording remains unchanged.
\SourcePage{121}

\providecommand{\TohokuContentsChapter}[1]{%
  \par\medskip\noindent\textbf{#1}\par}
\providecommand{\TohokuContentsEntry}[2]{%
  \par\noindent\hangindent=2.8em\hangafter=1%
  \makebox[2.8em][l]{#1}#2\par}

\LayerTarget{tohoku:introduction:III:p1}
des faisceaux, un exposé satisfaisant se trouvera dans [9]. De nombreuses
autres lacunes se signaleront d'elles mêmes à l'attention du lecteur.

\LayerTarget{tohoku:introduction:III:p2}
Pour terminer, je suis heureux d'exprimer mes remerciements à MM.
R. Godement, H. Cartan et J. P. Serre, dont l'intérêt a été le stimulant
indispensable pour la rédaction du présent travail.

\LayerTarget{tohoku:contents}
\LayerContents{Table des matières}
\begin{center}
  \bfseries TABLE DES MATIÈRES
\end{center}

\begingroup\small
\LayerTarget{tohoku:chapter:I}
\TohokuContentsChapter{CHAPITRE I\quad GÉNÉRALITÉS SUR LES CATÉGORIES
ABÉLIENNES}

\LayerTarget{tohoku:section:1.1}
\TohokuContentsEntry{1.1.}{Catégories}
\LayerTarget{tohoku:section:1.2}
\TohokuContentsEntry{1.2.}{Foncteurs}
\LayerTarget{tohoku:section:1.3}
\TohokuContentsEntry{1.3.}{Catégories additives}
\LayerTarget{tohoku:section:1.4}
\TohokuContentsEntry{1.4.}{Catégories abéliennes}
\LayerTarget{tohoku:section:1.5}
\TohokuContentsEntry{1.5.}{Sommes et produits infinis}
\LayerTarget{tohoku:section:1.6}
\TohokuContentsEntry{1.6.}{Catégories de diagrammes, propriétés de permanence}
\LayerTarget{tohoku:section:1.7}
\TohokuContentsEntry{1.7.}{Exemples de catégories définies par des schémas de
diagrammes}
\LayerTarget{tohoku:section:1.8}
\TohokuContentsEntry{1.8.}{Limites inductives et projectives}
\LayerTarget{tohoku:section:1.9}
\TohokuContentsEntry{1.9.}{Générateurs et cogénérateurs}
\LayerTarget{tohoku:section:1.10}
\TohokuContentsEntry{1.10.}{Objets injectifs et projectifs.}
\LayerTarget{tohoku:section:1.11}
\TohokuContentsEntry{1.11.}{Catégories quotient}

\LayerTarget{tohoku:chapter:II}
\TohokuContentsChapter{CHAPITRE II\quad ALGÈBRE HOMOLOGIQUE DANS LES
CATÉGORIES ABÉLIENNES}

\LayerTarget{tohoku:section:2.1}
\TohokuContentsEntry{2.1.}{$\partial$-foncteurs et
$\partial^{*}$-foncteurs}
\LayerTarget{tohoku:section:2.2}
\TohokuContentsEntry{2.2.}{$\partial$-foncteurs dérivés}
\LayerTarget{tohoku:section:2.3}
\TohokuContentsEntry{2.3.}{Foncteurs dérivés}
\LayerTarget{tohoku:section:2.4}
\TohokuContentsEntry{2.4.}{Suites spectrales et foncteurs spectraux}
\LayerTarget{tohoku:section:2.5}
\TohokuContentsEntry{2.5.}{Foncteurs résolvants}

\LayerTarget{tohoku:chapter:III}
\TohokuContentsChapter{CHAPITRE III\quad COHOMOLOGIE A COEFFICIENTS DANS UN
FAISCEAU}

\LayerTarget{tohoku:section:3.1}
\TohokuContentsEntry{3.1.}{Généralités sur les faisceaux}
\LayerTarget{tohoku:section:3.2}
\TohokuContentsEntry{3.2.}{Définition des $H_{\Phi}^{p}(X,F)$}
\LayerTarget{tohoku:section:3.3}
\TohokuContentsEntry{3.3.}{Critères d'acyclicité}
\LayerTarget{tohoku:section:3.4}
\TohokuContentsEntry{3.4.}{Applications à des questions de relèvement du
groupe structural}
\LayerTarget{tohoku:section:3.5}
\TohokuContentsEntry{3.5.}{La suite exacte relative à un sous-espace fermé}
\LayerTarget{tohoku:section:3.6}
\TohokuContentsEntry{3.6.}{Sur la dimension cohomologique de certains espaces}
\LayerTarget{tohoku:section:3.7}
\TohokuContentsEntry{3.7.}{La suite spectrale de Leray d'une application
continue}
\LayerTarget{tohoku:section:3.8}
\TohokuContentsEntry{3.8.}{Comparaison avec la cohomologie de Čech}
\LayerTarget{tohoku:section:3.9}
\TohokuContentsEntry{3.9.}{Critères d'acyclicité par la méthode des
recouvrements}
\LayerTarget{tohoku:section:3.10}
\TohokuContentsEntry{3.10.}{Passages à la limite en cohomologie des faisceaux}

\LayerTarget{tohoku:chapter:IV}
\TohokuContentsChapter{CHAPITRE IV\quad LES EXT DE FAISCEAUX DE MODULES}

\LayerTarget{tohoku:section:4.1}
\TohokuContentsEntry{4.1.}{Les foncteurs $\mathrm{Hom}_{\mathbf{O}}(A,B)$ et
$\mathbf{Hom}_{\mathbf{O}}(A,B)$}
\LayerTarget{tohoku:section:4.2}
\TohokuContentsEntry{4.2.}{Les foncteurs $\mathrm{Ext}_{\mathbf{O}}^{n}(X;\,A,B)$ et
$\mathbf{Ext}^{n}(A,B)$ et la suite spectrale}
\endgroup

\LayerTarget{tohoku:introduction:footnote:1bis}
\SourceFootnote{1 bis}{M. P. Cartier vient de trouver une
formulation satisfaisante générale pour les structures multiplicatives en
Algèbre Homologique, qu'il exposera en son lieu.}

\TohokuContinuation{122}
